% GENERATED FILE. Edit canonical lessons, then run npm run generate:books.
% canonical-source-hash: fnv1a64:8c552d94f3ebaaed
% canonical-lessons: ZH-C63-dianhua, ZH-W63-dian-electric, ZH-C63-dadianhua, ZH-W63-da-hit, ZH-C63-wei, ZH-C63-haoma, ZH-C63-zhao, ZH-R63-review-58-63

\chapter{On the Phone}
\label{ch:zh-on-the-phone}
\hlchaptermodality{63}

\begin{chapteropening}
\textbf{By the end of this chapter:} \emph{I can answer the phone, ask for a number and say what I am looking for.}

Complete the last lesson of chapter 63: I can answer the phone, ask for a number and say what I am looking for.
\end{chapteropening}

\section[\texorpdfstring{\zh{电话} (diànhuà)}{电话 (diànhuà)}]{\zh{电话} (diànhuà) — telephone; phone call}

\label{lesson:ZH-C63-dianhua}



\begin{quote}
Before the new one: say the Mandarin for near, then the Mandarin for to know a fact.
\end{quote}

\subsection*{You'll want to know: \zh{电话}}
\textbf{\zh{电话}} — \emph{diànhuà} — \textquotedblleft{}telephone; phone call\textquotedblright{}.

\begin{quote}
\textbf{\zh{你有电话吗}} \emph{nǐ yǒu diànhuà ma} — Have you got a phone?
\end{quote}

The new character, in pieces: \zh{电} — \zh{日} with one long stroke running down through its middle, bending right along the bottom and hooking up. The next lesson puts it in your hand.

\begin{cousinweb}
\textbf{\zh{电话}} (\emph{diànhuà}) is built from \zh{电} 'lightning', now 'electricity', and \zh{话} 'talk', both inherited from Middle Chinese, \emph{denH} and \emph{hwaejH} (Guangyun, in Baxter's transcription): electric talk.
\end{cousinweb}

\begin{grammarlens}[title={What you've built}]
One more word for asking politely.
\end{grammarlens}

\subsection*{Your turn}
\begin{itemize}
\raggedright
  \item \emph{Say it:} \emph{diànhuà}
  \item \emph{Say it:} \emph{diànhuà}, once more
  \item \emph{Say it:} \emph{nǐ yǒu diànhuà ma} — Have you got a phone?
  \item \emph{Recall:} say the Mandarin for near, then the Mandarin for to know a fact, then say \emph{diànhuà} again
\end{itemize}

\begin{tcolorbox}[breakable,colback=teal!4,colframe=teal!35!black,title={Before you move on}]
What does \zh{电话} mean? (\textquotedblleft{}Telephone; phone call\textquotedblright{}.) Say it once more.
\end{tcolorbox}
\section[\texorpdfstring{\zh{电} (diàn)}{电 (diàn)}]{\zh{电} — \zh{日} with a long tail}

\label{lesson:ZH-W63-dian-electric}



\begin{quote}
The word first, then the character inside it.

\begin{itemize}
\raggedright
  \item \emph{Recall:} say \emph{diànhuà} — \textbf{R1}, one lesson back, before it decays
\end{itemize}

\begin{quote}
\zh{日} \zh{话}
\end{quote}

\emph{Read it:} each one aloud, with nothing said first — \emph{rì}, \emph{huà}
\end{quote}

\subsection*{Script — \zh{日} with a long tail}
\hlblockfigure{\detokenize{figures/ZH-W63-dian-electric-filmstrip.pdf}}{How \zh{电} is written, stroke by stroke}

\begin{quote}
\zh{电}
\end{quote}

\textbf{Two parts: the box first, then the tail.}

\textbf{\zh{日}, four strokes}: the left side, the top that turns down the right side, the middle bar, and the bottom bar.

\textbf{The tail, one stroke}: from above the top of the box, a vertical straight down through the middle, bending right along the bottom and hooking up at the end.

4 + 1 = five.

\subsection*{Your turn}
\begin{itemize}
\raggedright
  \item \emph{Write it:} \zh{电}, all five, with the model in front of you
  \item \emph{Cover:} write it again from memory, and count the pieces, not the strokes
\end{itemize}

\begin{tcolorbox}[breakable,colback=teal!4,colframe=teal!35!black,title={Before you move on}]
How many strokes in \zh{电}? (\textbf{Five}.) Which character you know is inside it? (\textbf{\zh{日}}.) Where does the last stroke start? (\textbf{Above the box, and it runs down through the middle}.)

Source: \href{https://raw.githubusercontent.com/chanind/hanzi-writer-data/68d10a4b21150cae5e1ebbd223eed289cf32d90c/data/\%E7\%94\%B5.json}{Hanzi Writer Data \zh{电}, pinned stroke paths and medians}.
\end{tcolorbox}
\section[\texorpdfstring{\zh{打电话} (dǎ diànhuà)}{打电话 (dǎ diànhuà)}]{\zh{打电话} (dǎ diànhuà) — to phone; to make a phone call}

\label{lesson:ZH-C63-dadianhua}



\begin{quote}
Before the new one: say the Mandarin for to know a fact, then the Mandarin for telephone.
\end{quote}

\subsection*{You'll want to know: \zh{打电话}}
\textbf{\zh{打电话}} — \emph{dǎ diànhuà} — \textquotedblleft{}to phone; to make a phone call\textquotedblright{}.

\begin{quote}
\textbf{\zh{我想打电话}} \emph{wǒ xiǎng dǎ diànhuà} — I'd like to make a phone call.
\end{quote}

The new character, in pieces: \zh{打} — the hand piece, a narrow \zh{手}, on the left, and on the right a two-stroke piece read dīng. The next lesson puts it in your hand.

\begin{cousinweb}
\textbf{\zh{打电话}} (\emph{dǎ diànhuà}) is built from dǎ 'to strike' and \zh{电话} 'telephone'. The first is inherited from Middle Chinese \emph{taengX} (Guangyun, in Baxter's transcription), glossed 'to strike', though dǎ is not the sound \emph{taengX} would regularly give. Making a call is not among its senses there, so how the verb came to do it is not settled.
\end{cousinweb}

\begin{grammarlens}[title={What you've built}]
One more word for asking politely.
\end{grammarlens}

\subsection*{Your turn}
\begin{itemize}
\raggedright
  \item \emph{Say it:} \emph{dǎ diànhuà}
  \item \emph{Say it:} \emph{dǎ diànhuà}, once more
  \item \emph{Say it:} \emph{wǒ xiǎng dǎ diànhuà} — I'd like to make a phone call.
  \item \emph{Recall:} say the Mandarin for to know a fact, then the Mandarin for telephone, then say \emph{dǎ diànhuà} again
\end{itemize}

\begin{tcolorbox}[breakable,colback=teal!4,colframe=teal!35!black,title={Before you move on}]
What does \zh{打电话} mean? (\textquotedblleft{}To phone; to make a phone call\textquotedblright{}.) Say it once more.
\end{tcolorbox}
\section[\texorpdfstring{\zh{打} (dǎ)}{打 (dǎ)}]{\zh{打} — the hand piece, then dīng}

\label{lesson:ZH-W63-da-hit}



\begin{quote}
The word first, then the character inside it.

\begin{itemize}
\raggedright
  \item \emph{Recall:} say \emph{dǎ diànhuà} — \textbf{R1}, one lesson back, before it decays
\end{itemize}

\begin{quote}
\zh{手} \zh{电}
\end{quote}

\emph{Read it:} each one aloud, with nothing said first — \emph{shǒu}, \emph{diàn}
\end{quote}

\subsection*{Script — the hand piece, then dīng}
\hlblockfigure{\detokenize{figures/ZH-W63-da-hit-filmstrip.pdf}}{How \zh{打} is written, stroke by stroke}

\begin{quote}
\zh{打}
\end{quote}

\textbf{Two pieces, left before right.}

\textbf{The hand piece, three strokes}, a narrow form of \zh{手}: a short horizontal, a vertical through it that hooks to the left at the foot, then a stroke that rises from the lower left up across the vertical.

\textbf{The right piece, two strokes, read dīng}: a horizontal across the top, then a vertical down from its middle that hooks to the left.

3 + 2 = five.

\subsection*{Your turn}
\begin{itemize}
\raggedright
  \item \emph{Write it:} \zh{打}, all five, with the model in front of you
  \item \emph{Cover:} write it again from memory, and count the pieces, not the strokes
\end{itemize}

\begin{tcolorbox}[breakable,colback=teal!4,colframe=teal!35!black,title={Before you move on}]
How many pieces in \zh{打}? (\textbf{Two}.) Which character you know is the hand piece a narrow form of? (\textbf{\zh{手}}.) Which way does the rising stroke run? (\textbf{From the lower left up to the right}.)

Source: \href{https://raw.githubusercontent.com/chanind/hanzi-writer-data/68d10a4b21150cae5e1ebbd223eed289cf32d90c/data/\%E6\%89\%93.json}{Hanzi Writer Data \zh{打}, pinned stroke paths and medians}.
\end{tcolorbox}
\section[wéi]{wéi — hello on the phone}

\label{lesson:ZH-C63-wei}



\begin{quote}
Before the new one: say the Mandarin for telephone, then the Mandarin for to phone.
\end{quote}

\subsection*{You'll want to know: wéi}
\emph{wéi} — \textquotedblleft{}hello on the phone\textquotedblright{}.

\begin{quote}
\emph{wéi nǐ tīng wǒ shuō} — Hello? Listen to me.
\end{quote}

Its characters come later; for now it is a word for the ear.

On the phone it is said with a rising tone, as a question: wéi?

\begin{cousinweb}
\emph{wéi} is the call you answer the phone with. No dictionary we use records an older form, so its history is not settled.
\end{cousinweb}

\begin{grammarlens}[title={What you've built}]
One more word for asking politely.
\end{grammarlens}

\subsection*{Your turn}
\begin{itemize}
\raggedright
  \item \emph{Say it:} \emph{wéi}
  \item \emph{Say it:} \emph{wéi}, once more
  \item \emph{Say it:} \emph{wéi nǐ tīng wǒ shuō} — Hello? Listen to me.
  \item \emph{Recall:} say the Mandarin for telephone, then the Mandarin for to phone, then say \emph{wéi} again
\end{itemize}

\begin{tcolorbox}[breakable,colback=teal!4,colframe=teal!35!black,title={Before you move on}]
What does wéi mean? (\textquotedblleft{}Hello on the phone\textquotedblright{}.) Say it once more.
\end{tcolorbox}
\section[hàomǎ]{hàomǎ — number (of a phone, a room); identifying number}

\label{lesson:ZH-C63-haoma}



\begin{quote}
Before the new one: say the Mandarin for to phone, then the Mandarin for hello on the phone.
\end{quote}

\subsection*{You'll want to know: hàomǎ}
\emph{hàomǎ} — \textquotedblleft{}number (of a phone, a room); identifying number\textquotedblright{}.

\begin{quote}
\emph{nǐ de diànhuà hàomǎ shì duōshao} — What's your phone number?
\end{quote}

Its characters come later; for now it is a word for the ear.

\begin{cousinweb}
\emph{hàomǎ} is built from hào and mǎ, both inherited from Middle Chinese, \emph{hawH} and \emph{maeX} (Guangyun, in Baxter's transcription). The dictionary glosses the first as a call or a command and the second as agate, a stone like jade; how the pair came to mean a number is not settled.
\end{cousinweb}

\begin{grammarlens}[title={What you've built}]
One more word for asking politely.
\end{grammarlens}

\subsection*{Your turn}
\begin{itemize}
\raggedright
  \item \emph{Say it:} \emph{hàomǎ}
  \item \emph{Say it:} \emph{hàomǎ}, once more
  \item \emph{Say it:} \emph{nǐ de diànhuà hàomǎ shì duōshao} — What's your phone number?
  \item \emph{Recall:} say the Mandarin for to phone, then the Mandarin for hello on the phone, then say \emph{hàomǎ} again
\end{itemize}

\begin{tcolorbox}[breakable,colback=teal!4,colframe=teal!35!black,title={Before you move on}]
What does hàomǎ mean? (\textquotedblleft{}Number (of a phone, a room); identifying number\textquotedblright{}.) Say it once more.
\end{tcolorbox}
\section[zhǎo]{zhǎo — to look for}

\label{lesson:ZH-C63-zhao}



\begin{quote}
Before the new one: say the Mandarin for hello on the phone, then the Mandarin for identifying number.
\end{quote}

\subsection*{You'll want to know: zhǎo}
\emph{zhǎo} — \textquotedblleft{}to look for\textquotedblright{}.

\begin{quote}
\emph{nǐ zhǎo shénme} — What are you looking for?
\end{quote}

Its characters come later; for now it is a word for the ear.

\begin{cousinweb}
\emph{zhǎo} is the everyday verb for looking for something. No dictionary we use records an older form, so its history is not settled.
\end{cousinweb}

\begin{grammarlens}[title={What you've built}]
One more word for asking politely.
\end{grammarlens}

\subsection*{Your turn}
\begin{itemize}
\raggedright
  \item \emph{Say it:} \emph{zhǎo}
  \item \emph{Say it:} \emph{zhǎo}, once more
  \item \emph{Say it:} \emph{nǐ zhǎo shénme} — What are you looking for?
  \item \emph{Recall:} say the Mandarin for hello on the phone, then the Mandarin for identifying number, then say \emph{zhǎo} again
\end{itemize}

\begin{tcolorbox}[breakable,colback=teal!4,colframe=teal!35!black,title={Before you move on}]
What does zhǎo mean? (\textquotedblleft{}To look for\textquotedblright{}.) Say it once more.
\end{tcolorbox}
\section[(dialogue)]{Looking back: from Things You Carry to On the Phone}

\label{lesson:ZH-R63-review-58-63}



\begin{quote}
Say the Mandarin for mobile phone and bag.
\end{quote}

\subsection*{Your turn}
\begin{itemize}
\raggedright
  \item \emph{Say it:} mobile phone, bag, computer, cup, clothes
  \item \emph{Say it:} to wear, shoe, hat, trousers, pretty
  \item \emph{Say it:} white, colour, red, black, blue
  \item \emph{Say it:} long, short in length, tall, short in height, new
  \item \emph{Say it:} fast, no longer new, far, near, to know a fact
  \item \emph{Say it:} telephone, to phone, hello on the phone, identifying number, to look for
  \item \emph{Read it:} \zh{机} \zh{白} \zh{长} \zh{快} \zh{电} \zh{打} — each one aloud, then a word that holds it
\end{itemize}

\begin{tcolorbox}[breakable,colback=teal!4,colframe=teal!35!black,title={Before you move on}]
Mobile phone? (\textbf{\zh{手机}}, \emph{shǒujī}.) Bag? (\emph{bāo}.) Computer? (\emph{diànnǎo}.) Cup? (\emph{bēizi}.) Clothes? (\emph{yīfu}.) To wear? (\emph{chuān}.) Shoe? (\emph{xié}.) Hat? (\emph{màozi}.) Trousers? (\emph{kùzi}.) Pretty? (\emph{piàoliang}.) White? (\textbf{\zh{白}}, \emph{bái}.) Colour? (\emph{yánsè}.) Red? (\emph{hóng}.) Black? (\emph{hēi}.) Blue? (\emph{lán}.) Long? (\textbf{\zh{长}}, \emph{cháng}.) Short in length? (\emph{duǎn}.) Tall? (\emph{gāo}.) Short in height? (\emph{ǎi}.) New? (\emph{xīn}.) Fast? (\textbf{\zh{快}}, \emph{kuài}.) No longer new? (\emph{jiù}.) Far? (\emph{yuǎn}.) Near? (\emph{jìn}.) To know a fact? (\emph{zhīdào}.) Telephone? (\textbf{\zh{电话}}, \emph{diànhuà}.) To phone? (\textbf{\zh{打电话}}, \emph{dǎ diànhuà}.) Hello on the phone? (\emph{wéi}.) Identifying number? (\emph{hàomǎ}.) To look for? (\emph{zhǎo}.)
\end{tcolorbox}
